% !TeX root = ../../Book4.tex
% 纲要标题：### 1.4 自然同构与典范同构
% 纲要位置：计划纲要.md，第 2469 行。

\section{自然同构与典范同构}
\label{b4:sec:0104}

同构的存在与同构的自然性是两件事。张量积、对偶及模分解中的映射，能说明哪些对应由结构决定，哪些对应仍取决于选基或选补模。

% 纲要内容：
%
% * 双对偶与张量积的自然同构
% * 选基同构与典范同构
% * 挠子模、无挠商与有限生成模的非自然分裂
%
% ---
%

\subsection{双对偶与张量积的自然同构}
\label{b4:subsec:0104:01}

第二卷已经在模的具体构造中使用自然性。本节用函子和自然变换描述这些映射，辨认其源范畴、目标范畴和分量。

\ProseParagraphBreak
自然性最初也来自一个具体困难：一族对象逐个同构，是否就能在取极限后得到同构？Eilenberg 与 Mac Lane 的一般理论起于研究同调中的这种问题。\footnote{取极限的研究动机见\cite[Introduction, p. 236, especially footnote 4]{EilenbergMacLane1945NaturalEquivalences}；该脚注说明，从复形同调转向空间同调时需要这种极限过程。}逐项同构还须与系统中的过渡映射交换，才是两套系统之间的同构；这种相容性正是自然性。他们以有限维实向量空间为例，对比选基得到的 \(V\cong V^*\) 与求值映射给出的 \(V\cong V^{**}\)：前者随基改变，后者同时与所有线性映射相容。\footnote{对偶与双对偶的例子见\cite[Introduction, pp. 232--234]{EilenbergMacLane1945NaturalEquivalences}。}下面把求值映射的自然性与有限维时的可逆性分别写清楚。

\ProseParagraphBreak
设 \(k\) 为域。两次对偶给出协变函子
\(D^2:\mathbf{Vect}_k\rightarrow\mathbf{Vect}_k\)，其中
\(D^2(V)=V^{**}\)、\(D^2(f)=f^{**}\)。求值映射
\[
j_V:V\longrightarrow V^{**},\qquad j_V(v)(\lambda)=\lambda(v)
\]
组成 \(\operatorname{Id}\Rightarrow D^2\) 的自然变换。第二卷定理1.2.9“有限维双对偶同构”已经证明
\[
f^{**}j_V=j_Wf
\]
对任意线性映射 \(f:V\rightarrow W\) 成立，并证明有限维时 \(j_V\) 可逆。结合\cref{b4:prop:natural-iso-components}，它在
\(\mathbf{Vect}_k^{\mathrm{fd}}\) 上是自然同构。有限维保证分量的满射性，自然性本身并不需要这个条件。

\[
\begin{tikzcd}[column sep=large,row sep=large]
V\arrow[r,"j_V"]\arrow[d,"f"']&V^{**}\arrow[d,"f^{**}"]\\
W\arrow[r,"j_W"']&W^{**}.
\end{tikzcd}
\]

再设 \(R\) 为交换含幺环，\(L,M,N\) 为幺性 \(R\) 模。第二卷定理6.2.3“结合与单位同构”已经构造
\[
\begin{aligned}
\lambda_M &:R\otimes_RM\longrightarrow M,
&r\otimes m&\longmapsto rm,\\
\alpha_{L,M,N}&:(L\otimes_RM)\otimes_RN
\longrightarrow L\otimes_R(M\otimes_RN),
&(l\otimes m)\otimes n&\longmapsto l\otimes(m\otimes n).
\end{aligned}
\]
第一族是两个从 \(R\text{-}\mathbf{Mod}\) 到自身的函子之间的自然同构，第二族则从三个模范畴的积范畴出发。若分别给三个变量施以同态 \(f,g,h\)，两种复合在纯张量上的值都是
\(f(l)\otimes(g(m)\otimes h(n))\)，因此三变量自然性有明确的可核验含义。此处采用交换环以免另记外侧双模作用；一般双模版本仍以第二卷的陈述为准。

\ProseParagraphBreak
这两个例子也说明一个同构族必须说清楚相对于哪些态射自然。双对偶涉及所有线性映射，张量结合涉及三个变量上的同态。只说“公式没有写基”可以帮助猜测自然性，却不能代替相应方块的等式。

\subsection{选基同构与典范同构}
\label{b4:subsec:0104:02}

通常把由给定数据直接确定、或由明确的相容条件唯一刻画的映射称为典范映射。“典范”需要指出相对于哪份数据；自然性则是给定两个函子以后可以检验的一组等式。两者有密切联系，但不是同一个定义。

\begin{proposition}\label{b4:prop:natural-scalars}
设 \(k\) 为域。对向量空间及全部线性映射组成的范畴 \(\mathbf{Vect}_k\)，每个自然变换
\(\eta:\operatorname{Id}_{\mathbf{Vect}_k}\Rightarrow\operatorname{Id}_{\mathbf{Vect}_k}\)
都唯一具有形式
\[
\eta_V=c\,1_V\qquad(V\in\mathbf{Vect}_k),
\]
其中 \(c\in k\)。它是自然同构当且仅当 \(c\ne0\)。
\end{proposition}
\begin{proof}
令 \(c=\eta_k(1)\)。由于 \(\eta_k\) 线性，它为乘 \(c\) 的映射。任取 \(v\in V\)，线性映射 \(a_v:k\rightarrow V\)、\(r\mapsto rv\)，使自然性给出
\[
\eta_V(v)=\eta_Va_v(1)=a_v\eta_k(1)=cv.
\]
所以 \(c\) 决定全部分量，并由 \(V=k\) 保证唯一。反过来，标量乘法与每个线性映射交换，所以确实自然。各分量可逆当且仅当在 \(k\) 上可逆，也就是 \(c\ne0\)。
\end{proof}

因此即使已经要求自然，同构族也未必唯一：每个非零 \(c\) 都给出一个自然自同构。若另要求在 \(k\) 上把 \(1\) 送到 \(1\)，才得到恒等族。使用“唯一的典范同构”时，必须说出这样的确定条件，或给出相关泛性质。

\begin{example}\label{b4:ex:basis-dual-not-natural}
在实直线 \(V=\mathbb R e\) 中，以基 \(e\) 规定同构
\(b_e:V\rightarrow V^*\)、\(e\mapsto e^*\)。改用基 \(e'=2e\)，则其对偶基为
\((e')^*=\dfrac12e^*\)，所以
\[
b_{e'}(e)=\dfrac14e^*,\qquad b_e(e)=e^*.
\]
选基产生了同构，却没有产生同一个映射。

\ProseParagraphBreak
更严格地说，普通对偶反变，不能把 \(\operatorname{Id}\) 和 \(D\) 直接作为两个同向函子。令 \(\mathcal V_{\mathrm{iso}}\) 为有限维实向量空间及其间线性同构组成的范畴。在这里，对偶可以按下列协变规则重新组织：
\[
D_{\mathrm{iso}}:\mathcal V_{\mathrm{iso}}\longrightarrow\mathcal V_{\mathrm{iso}},
\qquad V\longmapsto V^*,\quad f\longmapsto(f^{-1})^*.
\]
假如存在自然同构
\(b:\operatorname{Id}_{\mathcal V_{\mathrm{iso}}}\Rightarrow D_{\mathrm{iso}}\)，对非零有限维实空间的同构
\(f=2\,1_V\)，自然性将要求
\[
\dfrac12 b_V=b_V(2\,1_V)=2b_V,
\]
故 \(b_V=0\)，矛盾。这个反例针对的是这两个确定的函子及所有实线性同构。
\end{example}

不过，不能从这个例子推出“使用选择的同构都不自然”。如果连同函子对态射的规则也一起改变，选定的同构反而可以组成自然同构。下面写清这一步。

\begin{proposition}\label{b4:prop:transported-functor}
设 \(\mathcal C,\mathcal D\) 为范畴，\(F:\mathcal C\rightarrow\mathcal D\) 为函子。假定对每个 \(X\in\mathcal C\) 已指定对象 \(G_X\in\mathcal D\) 及同构
\(\theta_X:F(X)\rightarrow G_X\)。则唯一存在函子
\(G:\mathcal C\rightarrow\mathcal D\)，其对象值为 \(G_X\)，且这些 \(\theta_X\) 组成自然同构 \(F\Rightarrow G\)。它的态射作用为
\[
G(f)=\theta_YF(f)\theta_X^{-1}
\qquad(f:X\rightarrow Y).
\]
\end{proposition}
\begin{proof}
自然性 \(G(f)\theta_X=\theta_YF(f)\) 已经强制给出所列公式，所以若存在就唯一。按公式定义后，恒等箭头映为恒等；对可复合 \(f,g\)，中间的
\(\theta_Y^{-1}\theta_Y\) 消去，得到 \(G(g)G(f)=G(gf)\)。公式本身又给出自然性，分量可逆保证自然同构。
\end{proof}

例如，设 \(k\) 为域，\(\mathcal V\) 是由一族有限维 \(k\) 向量空间及其间全部线性映射组成的全子范畴。记包含函子为
\(J:\mathcal V\hookrightarrow\mathbf{Vect}_k^{\mathrm{fd}}\)，并为每个 \(V\in\mathcal V\) 指定一组基，得到坐标同构
\(\theta_V:V\rightarrow k^{\dim V}\)。上一命题给出函子
\[
\begin{aligned}
G&:\mathcal V\longrightarrow\mathbf{Vect}_k^{\mathrm{fd}},\\
G(V)&=k^{\dim V},\qquad
G(f)=\theta_Wf\theta_V^{-1}\quad(f:V\longrightarrow W).
\end{aligned}
\]
其态射作用就是相应的矩阵规则，而 \(\theta:J\Rightarrow G\) 是自然同构。标准坐标空间不必属于 \(\mathcal V\)，所以这里比较的是两个取值于 \(\mathbf{Vect}_k^{\mathrm{fd}}\) 的函子。所选基没有消失，而是进入了 \(G\) 的定义；这些基族作为数据给定，其大小约定见\subsecref{b4:subsec:0101:03}。

\subsection{模结构定理中的自然性}
\label{b4:subsec:0104:03}

第二卷定理4.3.2已经证明：主理想整环 \(R\) 上的有限生成模具有不变因子分解
\[
M\cong R^r\oplus R/(d_1)\oplus\cdots\oplus R/(d_s),
\qquad d_1\mid\cdots\mid d_s,
\]
其中 \(d_i\) 为非零非单位。定理4.3.4进一步证明，自由秩和不变因子的相伴类由 \(M\) 的同构类型唯一确定。这并没有为每个 \(M\) 指定一个唯一分解同构；要讨论分解随同态怎样变化，还须分别考察其中的子模、商模与分裂映射。

\begin{proposition}\label{b4:prop:torsion-quotient-functor}
设 \(R\) 为交换整环。对每个幺性左 \(R\) 模 \(M\)，令
\[
t_R(M)=\{m\in M:\text{存在 }0\ne a\in R,\ am=0\},
\qquad Q_R(M)=M/t_R(M).
\]
则 \(t_R,Q_R\) 都是 \(R\text{-}\mathbf{Mod}\) 上的函子。包含
\(t_R(M)\hookrightarrow M\) 与商映射 \(M\rightarrow Q_R(M)\) 都随 \(M\) 自然。
\end{proposition}
\begin{proof}
若 \(am=0\)、\(bn=0\)，其中 \(a,b\ne0\)，则
\(ab(m-n)=0\)，且 \(ab\ne0\)；标量倍也仍被 \(a\) 零化，所以 \(t_R(M)\) 是子模。对同态 \(f:M\rightarrow N\)，有
\(a\cdot f(m)=f(am)=0\)，故 \(f(t_R(M))\subseteq t_R(N)\)。限制与诱导商映射于是良定义；它们对恒等与复合的相容性直接由限制和商的定义得到。包含方块显然交换；商方块在元素 \(m+t_R(M)\) 上的两个像都为 \(f(m)+t_R(N)\)。
\end{proof}

当 \(R\) 为 PID、\(M\) 有限生成时，已证的结构定理说明 \(Q_R(M)\) 有限自由，短正合列
\[
0\longrightarrow t_R(M)\longrightarrow M
\longrightarrow Q_R(M)\longrightarrow0
\]
可以分裂。但选择一个分裂比确定挠子模和商模要求更多，它一般不能对全部同态自然地进行。

\begin{proposition}\label{b4:prop:no-natural-torsion-section}
对每个有限生成交换群 \(A\)，记 \(t(A)\) 为其挠子群、\(q_A:A\rightarrow A/t(A)\) 为商映射。在有限生成交换群及群同态的范畴中，不存在一族群同态截面
\[
s_A:A/t(A)\longrightarrow A,\qquad q_As_A=1_{A/t(A)},
\]
使 \(s\) 对每个群同态自然。
\end{proposition}
\begin{proof}
只需考察 \(A=\mathbb Z\oplus\mathbb Z/3\mathbb Z\)。把商识别为 \(\mathbb Z\)，任一截面满足
\(s_A(1)=(1,\bar c)\)。自同构
\[
u:A\longrightarrow A,\qquad u(n,\bar a)=(n,\bar a+\bar n)
\]
在商 \(A/t(A)\) 上诱导恒等。若截面自然，就有 \(us_A=s_A\)。在 \(1\) 处比较得到
\((1,\bar c+\bar1)=(1,\bar c)\)，矛盾。
\end{proof}

对这里讨论的每个有限生成交换群，其挠子群商序列都存在截面；命题否定的是把这些截面选成对全部群同态自然的一族。这也回答了章首提出的问题：各个群分别存在的分解，未必能够同时与所有群同态相容。自然的子对象和自然的商对象，并不自动给出一个自然直和分解。后续的正合性与导出函子需要记录的，正是这些映射和相容关系，而不仅是对象的同构类型。
